\section{A Key You Did Not Produce}
\label{sec:ch06-a-key-you-did-not-produce}

\index{key!arrives from outside|(}%
Every other resolver in this book is handed a value the schema has already
checked. An \code{Int!} argument is an integer by the time the method runs, a
\code{DateTime} has been coerced, and an input object was validated field by
field. The reference resolver is handed a JSON object typed \code{\_Any}, which
chapter~\ref{ch:federation-model} described as deliberately shapeless, and
shapeless means unchecked.

\index{representation!naming one type with another type key}%
So send the subgraph a representation that names a \code{Session} and carries a
\code{Speaker}'s node id. \code{U3BlYWtlcjox} is the node id
chapter~\ref{ch:schema-design} printed for the first speaker, and the operation
is section~\ref{sec:ch06-the-key-arrives-as-a-string}'s, unchanged:

\begin{minted}{text}
POST /graphql HTTP/1.1
Host: localhost:5001
Content-Type: application/json

{"query":"<the same operation>","variables":
{"r":[{"__typename":"Session","id":"U3BlYWtlcjox"}]}}
\end{minted}

\begin{minted}[fontsize=\footnotesize]{json}
{"data":{"_entities":[{"__typename":"Session","title":"Schemas That Outlive Their Authors"}]}}
\end{minted}

\index{entities@\code{\_entities}!answers the wrong object}%
HTTP 200. No \code{errors} key. The \code{\_\_typename} is the one the caller
asked for, the fragment matched, the field is populated, and every part of that
response is well formed. It is also a session nobody asked for.
\code{U3BlYWtlcjox} decodes to \code{Speaker:1}, and the only reason a session
came back at all is that speaker one and session one share the number~1.

\index{key!decodes to an integer whatever it names}%
The mechanism is duller than the result. \code{Parse} was asked for an
\code{int}, and there is an \code{int} in that string, so it returned one.
Nothing was malformed. The base64 decoded cleanly, the integer was real, and the
resolver spent it. A schema check cannot see this, composition cannot see it,
and the type system agrees with every line.

\index{key!both sides must agree, and one side must check}%
Chapter~\ref{ch:federation-model} said both sides of a seam have to compute the
same string for the same object, and that nothing checks it. This is the other
half of that sentence and it is the half you can act on. The side sending a key
cannot be made to send the right one. The side receiving one can look at what it
was given.

\subsection{Reading the Other Half}

\index{NodeId@\code{NodeId}!the type name was there all along}%
The information needed to refuse that request was in the resolver's hand the
whole time. \code{Parse} returned a \code{NodeId} carrying two parts, and
section~\ref{sec:ch06-the-key-arrives-as-a-string} spent one of them. Here is
the session's type class again, using both:

\begin{minted}[highlightlines={54-64}]{csharp}
using GreenDonut.Data;
using HotChocolate.ApolloFederation.Resolvers;
using HotChocolate.ApolloFederation.Types;
using HotChocolate.Types.Relay;
using Microsoft.EntityFrameworkCore;

namespace Sessions;

[ObjectType<Session>]
[Key("id")]
[ReferenceResolver(
    EntityResolver = nameof(ResolveSessionReferenceAsync))]
public static partial class SessionType
{
    /// <summary>
    /// The person giving this session.
    /// </summary>
    public static async Task<Speaker?> GetSpeakerAsync(
        [Parent(nameof(Session.SpeakerId))] Session session,
        ISpeakerByIdDataLoader speakerById,
        CancellationToken ct)
        => await speakerById.LoadAsync(session.SpeakerId, ct);

    /// <summary>
    /// Loads one session by the identifier inside a global node id.
    /// </summary>
    [NodeResolver]
    public static async Task<Session?> GetSessionAsync(
        int id,
        ConferenceContext db,
        QueryContext<Session> query,
        CancellationToken ct)
        => await db.Sessions
            .Where(s => s.Id == id)
            .With(query)
            .FirstOrDefaultAsync(ct);

    /// <summary>
    /// Answers one representation from the entity route. Two things
    /// differ from the node resolver above, and neither is a choice.
    /// The key arrives as the undecoded node id string, so this method
    /// decodes it. And a QueryContext cannot be injected here, so there
    /// is no projection and every column is selected. One statement per
    /// representation is what chapter 10 measures and then fixes.
    /// </summary>
    [GraphQLIgnore]
    public static async Task<Session?> ResolveSessionReferenceAsync(
        string id,
        INodeIdSerializer serializer,
        ConferenceContext db,
        CancellationToken ct)
    {
        var nodeId = serializer.Parse(id, typeof(int));

        // Nothing has checked what is inside this string, and the
        // router is not the only caller the entity route has. A
        // Speaker's node id decodes to a perfectly good integer;
        // without this line it would answer with the session that
        // happens to carry it.
        if (nodeId.TypeName != nameof(Session))
        {
            return null;
        }

        var key = (int)nodeId.InternalId!;

        return await db.Sessions
            .Where(s => s.Id == key)
            .FirstOrDefaultAsync(ct);
    }
}
\end{minted}

\index{SpeakerType@\code{SpeakerType}!takes the same guard}%
The same block goes in the same place in \code{SpeakerType.cs}, between the
\code{Parse} and the load, comparing against \code{nameof(Speaker)} instead. The
two nouns in the comment swap with it: there it is a \code{Session}'s node id
that would otherwise answer with a speaker. Appendix~\ref{app:source} carries
both files in their final state.

\index{entities@\code{\_entities}!null is the specification's answer}%
Null is the correct answer here rather than a convenient one. The specification
says entries in the returned list \enquote{can be null if no entity exists for a
provided representation}~\autocite{apollosubgraphspec}, and a representation
naming a \code{Session} with a \code{Speaker}'s key names no session. The same
request now answers \code{\{"data":\{"\_entities":[null]\}\}}, and it does so
without issuing a statement, because the resolver refuses before it reaches the
table.

\index{node@\code{node}!reads the type name already}%
Note who was doing this all along. \code{node(id: "U3BlYWtlcjox")} against the
same service answers with a \code{Speaker}, because the machinery
chapter~\ref{ch:schema-design} turned on reads the type name out of the id and
dispatches on it. Global object identification never had this problem. The
entity route walked past the part of the mechanism that solved it and took the
integer without the name.

\subsection{Which Is Not the Ordering Failure}

\index{representation order!is a different failure}%
Two chapters of this book are about a subgraph answering with the wrong object,
and they are not the same failure. Chapter~\ref{ch:representation-order} is
about order: every object in the list is correct and the list is in the wrong
sequence, so each answer is attached to somebody else's question. This one is
about the key: the list is in the right order and one of the objects in it was
never asked for.

\index{failure!invisible to every gate so far}%
What they share is the shape of the symptom, which is the shape most of
part~III has: a 200 carrying well-typed data, no error key, and nothing in any
log. Neither is caught by a schema check, by composition, by the compiler or by
a client. Both are caught by the resolver being written correctly the first
time, which is the argument for reading the key rather than trusting it, and
the reason I would rather spend two lines here than a morning in
chapter~\ref{ch:blame-routing}'s territory later.
\index{key!arrives from outside|)}%
